\subsection*{C19. Extended individual, team and cross-view results}

\paragraph{Individual geometric and labor responses.}
On 201 images shared by C1, C2 and C3, mean individual Q was 64.965, 66.913 and 80.971, respectively. Mean matched-boundary errors were $1.856^\circ$, $1.730^\circ$ and $0.941^\circ$. This error averages absolute upper- and lower-boundary angular deviations at matched corner pairs and differs from footprint boundary distance $S_b$. On the 130 easy images shared by C2 and C3, their turn-count difference was 0.058 while their Q difference was 13.943 points.

Within-person, within-building difficult-minus-easy comparisons gave C3 an additional 2.330 turns relative to C2 across fifteen buildings, positive in all fifteen. The corresponding Q difference was +5.235 points, positive in nine buildings. Adding geometric features to the difficulty, annotator and source controls reduced held-out-building turn-deficit MSE from 1.744 to 1.163. Reference turn count is among those predictors and also enters the definition of turn deficit; binary simplification MSE and Q MAE are reported alongside it in the main text. Supplement C7 specifies the prediction procedure.

In 100 slower--faster pairs reusing 35 annotators, the slower person's easy-to-difficult advantage changed by $-$0.054 turns and $-$0.029 short edges on average. Positive changes occurred in 53 and 55 pairs, respectively. Across 162 images and seventeen buildings, J2 minus J1 gave Q $-$0.133 and a geometric mean labor ratio of 0.885. These outcomes describe recorded person-time and structural expression.

The six expert-confirmed controlled Semi tasks yielded 156 opportunities from 26 annotators, with 25 unchanged submissions. The retained ordinary primary-target subset has 138 responses and the issue-specific outcome counts in C11. Three separate natural-initialization submissions for e9zR4mvMWw7-39 also retained their initial coordinates. The natural initialization had been checked by experts before the experiment. Supplement B20 shows the natural example. In the duplicated-pair example in B8, panorama Z6MFQCViBuw-10, A2 removed the added pair while retaining every other coordinate and A3 retained it.

\paragraph{Small groups, completion and retrospective stability.}
On the fixed 171-image panel, the consecutive changes in mean Q for sizes 3--6, 6--10, 10--20 and 20--30 were +4.74, +1.75, +1.31 and +0.50 after rounding. The corresponding P10 gains were +8.27, +3.73, +3.78 and +2.85; each difference was calculated before rounding the displayed endpoint levels. Actual failure rates rose by 1.894 percentage points from three to six people and fell by 0.416 points from twenty to thirty. Complete-and-evaluable rates changed by $-$1.857 and +0.420 points. An output can be produced but unscorable, accounting for the distinction between these rates.

Among 98 images with capacity 35, 97 had scorable outputs at both thirty and 35 people. On these 97 images, mean Q increased by 0.281 points, P10 by 2.736 and within-image standard deviation fell by 1.952. Complete capacity strata and the full-pool diagnostics are given in C10.

Retrospective stability requires the range of image-mean Q from a candidate size through the actual full pool to be at most a specified tolerance, at least $\max(5,\lceil0.15N\rceil)$ additional members to remain observable, and every size to have at least 64 scorable paths. With a two-point tolerance, onsets were detected on 163 of 206 ordinary images: 123/130 easy, 25/39 moderate and 15/37 difficult images. The remaining 43 had no detected onset. Median sizes among detected onsets were nine, sixteen and seventeen. Tolerances of one, two, three and five points yielded 145, 163, 170 and 184 detected onsets, respectively. The criterion uses the complete observed curve.

Early-entrant composition tracked later quality within the stored natural paths. Orders beginning with six C2+C3 entrants averaged about 1.78 Q points higher over sizes 6--20 than those beginning with C1+C2 entrants, with positive differences on 165 of 190 comparable images. These orders vary in both profiles and particular people. Separately calculated one-to-two and two-to-three paired changes averaged about +4.23 and +3.61 Q points on the 171-image panel, requiring both endpoints scorable for each contrast. At two people, either person's support meets the at-least-half rule.

\paragraph{Additional composition and replacement comparisons.}
Increasing the C3 count by one at the expense of C2 in fixed 20-person groups raised composition-mean Q by 1.192 points, positive on 27 of 35 images. This comparison averages member sets within a composition. On common quality/time support for twelve-person teams, C2-to-C3 gave +1.963 Q and a geometric labor ratio of 1.020 on 27 images; J3-to-J4 gave +2.701 Q and 0.897 on fourteen images.

The direct automatic 20-person, seven-person C2-to-C3 replacement increased Q by 4.484 points, positive on 77 of 84 images. On its forty images with joint quality and time support, the increase was 4.406 Q points and 121.05 person-seconds. For the 14-person A--B--C comparison, the automatic-only support comprises 67 images in twenty buildings: $G_S=3.767$, positive on 49 images. The main comparison including assisted ordering has 69 images, $G_S=3.441$ and fifty positive image effects.

\begin{table}[ht]
\centering\small
\caption{Bidirectional six-person comparisons replacing three people. Both directions require all four outputs scorable on the same 36 images in ten buildings.}
\label{tab:bidirectional-half}
\begin{tabular}{lrrr}
\toprule
Direction & Symmetric excess $G_S$ & Positive excess & $Q_C-Q_B$\\
\midrule
C1 to C3 & 3.269 & 32/36 & +7.327\\
C3 to C1 & 3.955 & 24/36 & $-$7.777\\
\bottomrule
\end{tabular}
\end{table}
On this shared support, the C3-versus-C1 team-type mean difference is +7.865 Q points. It compares team types across the same images, separately from the two directional exchanges in the table.

The R1/R2 consensus-response profiles also differed in individual geometry. Across 192 common images, R2 exceeded R1 by 13.937 Q points, positive on 175; its matched-boundary error was $0.849^\circ$ lower and its omitted-area fraction 0.0545 lower. In the twelve-person, four-person matched replacement, within-profile Q-change magnitude averaged 4.132 and cross-profile magnitude 6.308, with the latter larger on 62 of 83 images.

\paragraph{Retained backgrounds and geometry.}
In the joint-profile comparison holding seven-person blocks fixed, increasing total size from twenty to 28 reduced absolute Q change by 1.317 points on 76 images, with reductions on 63. All four outputs were scorable on each path; the mean difference in signed effects between sizes was +0.014 points.

The background-reversal comparison fixes image, reference, team size, replacement size and exchanged members, using unchanged automatic endpoints. Of 276 blocks observed under at least two backgrounds, 56 reversed sign across seven images; 24 crossed both +1 and $-$1 Q point across five images. Requiring identical background profile counts yielded 120 subgroups, representing 108 distinct blocks across seven images and four buildings. Eighteen subgroups reversed sign; nine crossed both thresholds across four images. Blocks have differing numbers of observed backgrounds. The q9vSo1VnCiC-30 example in the main text had the largest effect range among comparisons with the same exchanged members and background profile counts.

Score range and geometric components explain additional response differences. In the automatic 20-person, seven-person C2-to-C3 comparison, 1,157 pairs with both Q values at least ninety differed by 2.064 points on average; 672 pairs with at least one lower value differed by 11.012. An image can contribute pairs to both sets. In a separate eight-person, four-person C3-to-C1 comparison on jtcxE69GiFV-13, A, B and C scored 72.914, 82.144 and 72.322. B retained six selected nodes, while A and C retained four; complete connection in this example used source-order review.

Among the seven of 84 images whose Q fell while prism IoU rose in the automatic twenty-person C2-to-C3 comparison, x8F5xyUWy9e-03 changed by $-$4.726 Q points, +2.450 IoU percentage points and +$0.836^\circ$ direction error. The complete outputs and component values explain the opposed quality summaries.

\paragraph{Cross-view structure and quality.}
Across the 128 directly related view pairs, same-person versus different-person mean gaps were 0.671 versus 0.752 effective turns, 0.377 versus 0.445 turn deficits and 0.363 versus 0.414 short-edge counts. Both comparisons use the same common-worker pool within a pair and equal pair weights. Across all pairs, simplification agreement was 88.22\% for the same person and 86.37\% for different people. Uniformly unsimplified responses left many pairs with little variation; the fifteen pairs with short edges in both references yielded the larger agreement difference reported in the main text.

Across four views of the illustrated room, A40 retained 10, 14, 16 and 16 effective turns, matching the reference counts; A41 retained 8, 4, 9 and twelve. Among eighteen common annotators in the two illustrated views, the door-side endpoint was matched in both views by five people, neither by six and one by seven; the vanity-side counts were five, eight and five. The counts agree at azimuth tolerances of 3, 6 and 9 degrees. The second manually confirmed room correspondence yielded inner-endpoint status changes for 19/33, 12/30 and 14/30 common annotators across three view pairs; its same-person endpoint agreement did not consistently exceed the different-person comparison.

Same-room mean-Q gaps averaged 5.051 points, compared with 8.931 for controls matched on building, difficulty, reference turns and geometrically hidden reference vertices. The comparison used common workers and 215 directional comparisons involving 113 same-room pairs in nine buildings. Gaps were smaller in eight buildings. Control pairs had unknown recorded relation and partially verified physical-room identity.

Within-pair Q rank correlation was positive on 105 of 128 pairs; forty of 57 people repeated relative-performance direction on more than half their pairs. Raw individual cross-view Q correlation was 0.401, peer-centered correlation 0.199 and correlation after adjustment for building-external personal history 0.013. The image-mean prediction uses both directions of 128 pairs, giving 256 comparisons; C7 defines its target-image and building aggregation and the separate individual-Q and structural prediction estimands.

The full-composition transfer analysis uses 324 panels from 127 directly related ordinary image pairs in thirteen buildings. Moments are averaged by direct pair, declared room group and building; aggregated covariance and variances yield the pooled correlation of $-$0.052. The specified replacement instead follows the same retained and exchanged people across nineteen view pairs in seven buildings. C14 reports automatic and assisted branches and their corresponding event supports.

\paragraph{Doorway shared-path comparison.}
On the fixed seventeen doorway images, screening excludes 145 of 584 submissions and retains 439. At ten people, the screening-minus-unscreened mean Q difference is +1.972 on the 3,411 of 4,352 paths scorable in both conditions. The main conditional-mean difference of +1.581 uses each condition's own scorable outputs. Both retain the same whole-space reference and underlying base orders. The complete 3--34-person doorway analysis uses decreasing image support above fourteen, whereas the main 3--14 curves retain all seventeen images.
